% Research fragment.  Dependencies: model_kernel.tex, uniformity.tex,
% inert_comparison.tex and pressure.tex in main revision 332617a.
% No claim of irrationality follows from this fragment.

\section{The full-environment kernel has critical essential radius one}
\label{ess:section}

Write $N(x)$ for the number of real vertices in a newborn environment
$x\in\mathcal X$.  All operators in this section act on
$\ell^\infty(\mathcal X)$, with the supremum norm.  Let
\[
 u_m(s)=\inf_x A_m(s;x),\qquad
 v_m(s)=\sup_x A_m(s;x)=\|\mathcal K_s^m\|.
\]
The previously proved inert comparison gives
\[
 \gamma_s A_m^\infty(s)\le u_m(s)\le A_m^\infty(s),
 \qquad \gamma_s=[12(s+3)]^{-1},
\]
and hence $u_m(s)^{1/m}\longrightarrow r_\infty(s)$.

\begin{proposition}[Finite blocks do not contain the critical root]\label{ess:finite}
For a finite subset $E\subset\mathcal X$, the compression
$\mathcal K_{s,E}$ is nilpotent, and
\[
 \det(I-z\mathcal K_{s,E})=1 \qquad(z\in\mathbb C).
\]
For a finite subset $F$ of the one-stub inert core states, the killed
proper-birth generator and the compressed core-ring kernel satisfy
\[
 \det(sI-L_F)
 =\prod_{y\in F}\left(s+N(y)-\frac1{D_0(y)}\right),
\]
and
\[
 \det(I-Q_{s,F})
 =\prod_{y\in F}\left(1-\frac1{D_0(y)(s+N(y))}\right)>0
 \quad(s>0).
\]
\end{proposition}
\begin{proof}
Every genealogical transition creates at least one real vertex and
retains the whole environment.  Ordering states by $N$ makes each
finite genealogical compression strictly upper triangular.

In the inert process, a proper birth increases $N$ by one, and its
total rate is $N-1/D_0$.  This proves the generator formula.  A
core-ring transition can retain the same population only if no leaf
birth preceded the ring and the root selected its stub.  Its discounted
mass is
\[
 Q_s(y,y)=\int_0^\infty e^{-(s+N(y))t}\frac{dt}{D_0(y)}
 =\frac1{D_0(y)(s+N(y))}.
\]
All other entries of the same population level vanish.  Thus $Q_{s,F}$
is upper triangular with the displayed diagonal.  Here $N\ge2$ and
$D_0\ge2$, so all factors are positive.
\end{proof}

\begin{theorem}[Essential-radius lower bound]\label{ess:radius}
For every $s>0$ and every integer $m\ge1$,
\[
 \|\mathcal K_s^m\|_{\rm ess}\ge u_m(s),
 \qquad
 r_{\rm ess}(\mathcal K_s)\ge r_\infty(s).
\]
Consequently
\[
 \boxed{r_{\rm ess}(\mathcal K_\mu)
       =r(\mathcal K_\mu)=1.}
\]
\end{theorem}
\begin{proof}
Fix $m,s$ and $\varepsilon>0$.  There are newborn environments of
arbitrarily large population.  Recursively choose such environments
$x_j$, and finite integers $M_j$, as follows.  All transitions counted
by $\mathcal K_s^m$ end in a finite environment of population at least
$N(x_j)+m$.  The corresponding row has finite mass
$A_m(s;x_j)$.  Hence monotone convergence permits a cutoff $M_j$ such
that
\[
 \mathcal K_s^m
 \mathbf1_{\{N(x_j)+m\le N\le M_j\}}(x_j)
 \ge A_m(s;x_j)-\varepsilon.
\]
Choose $x_{j+1}$ with $N(x_{j+1})>M_j$.  The resulting functions
\[
 f_j=\mathbf1_{\{N(x_j)+m\le N\le M_j\}}
\]
have pairwise disjoint supports and supremum norm at most one.

The sequence $(f_j)$ is weakly null in $\ell^\infty$.  Indeed, for
any bounded linear functional $\varphi$, choose scalars of modulus
one making each summand nonnegative; disjointness then gives
\[
 \sum_{j=1}^J|\varphi(f_j)|\le\|\varphi\|
 \quad\text{for every }J.
\]
In particular $\varphi(f_j)\to0$.  This argument uses disjoint bounded
population bands, not tail indicators; the latter need not be weakly
null in $\ell^\infty$.

If $S$ is compact, then $\|Sf_j\|\to0$.  Therefore
\[
 \|\mathcal K_s^m-S\|
 \ge (\mathcal K_s^mf_j)(x_j)-\|Sf_j\|
 \ge u_m(s)-\varepsilon-o(1).
\]
Taking the infimum over compact $S$ and then letting
$\varepsilon\downarrow0$ proves the essential-norm inequality.
The spectral-radius formula in the quotient by compact operators
gives
\[
 r_{\rm ess}(\mathcal K_s)
 =\lim_m\|\mathcal K_s^m\|_{\rm ess}^{1/m}
 \ge\lim_m u_m(s)^{1/m}=r_\infty(s).
\]

To identify the full spectral radius at the critical point, put
$R(s)=r(\mathcal K_s)=\lim_m v_m(s)^{1/m}$.  Each $A_m(s;x)$ is a
positive Laplace transform.  H\"older's inequality and the supremum
over $x$ show that $v_m(s)$ is log-convex.  The finite pointwise limit
\[
 \log R(s)=\lim_m\frac1m\log v_m(s)
\]
is therefore convex and continuous on $(0,\infty)$; finiteness follows
from
\[
 -\log[2(s+2)]\le\log R(s)\le-\log s.
\]
For $s>\mu$, the uniform newborn resolvent estimate gives
\[
 v_m(s)\le
 s\sup_x\int_0^\infty e^{-st}F_x(t)\,dt<\infty
 \quad\text{uniformly in }m.
\]
Thus $R(s)\le1$ for $s>\mu$, and continuity gives $R(\mu)\le1$.
On the other hand,
$R(\mu)\ge r_\infty(\mu)=1$.
Combining this with the essential lower bound proves the claim.
\end{proof}

\begin{corollary}[A second exact scalar characteristic]\label{ess:characteristic}
The full spectral radius $R(s)=r(\mathcal K_s)$ is continuous and
strictly decreasing on $(0,\infty)$ and
\[
 R(s)>1\ (s<\mu),\quad R(\mu)=1,\quad R(s)<1\ (s>\mu).
\]
\end{corollary}
\begin{proof}
The function $\log R$ is convex and nonincreasing.  If it were equal
at two distinct points, these two properties would force it to be
constant thereafter, contrary to $\log R(s)\le-\log s\to-\infty$.
\end{proof}

This theorem rules out a spectral-radius gap at the critical parameter
on the full supremum-norm space.  In particular, the full operator is not compact, and finite-rank
operators cannot approximate it in operator norm.  It does not rule
out a smaller invariant space of suitably local observables, or an
exact quotient that forgets remote population while retaining all
information needed for the growth exponent.  Nor does equality of
the spectral and essential spectral radii, by itself, assert that the
particular spectral point $1$ cannot be isolated.

Finite Schur complements of the displayed finite blocks cannot produce
a positive critical root.  Any useful scalar renewal reduction must
therefore introduce an exact quotient or retain a nontrivial contribution
from the boundary at infinite population.  A Fredholm or special-function
argument requires a proof for that new object; it cannot be imported
from the finite whole-environment compressions.
